\cleardoublepage
\thispagestyle{empty}

\begin{flushleft}
	\begin{tabular}{@{}l l@{}}
		\textbf{Title:} &
		\parbox[t]{0.65\textwidth}{\textit{The Structure of the World}}\\[0.4em]
		
		\textbf{Subtitle:} &
		\parbox[t]{0.65\textwidth}{Mathematics between Abstraction and Reality}\\[1.2em]
		
		\textbf{Author:} & Dipl.-Ing.\,(FH) Christian Weilharter\\[0.5em]
		\textbf{\copyright\ 2026} & Christian Weilharter, Traunstein\\[0.5em]
		\textbf{ISBN:} & 978-3-912302-13-4\\
		\textbf{Typesetting:} & \LaTeX\\
		\textbf{Printing:} & [Print-on-Demand Service]\\
		\textbf{Contact:} & info@mathandphysiks.eu\\
		\textbf{Website:} & www.mathandphysics.eu\\
	\end{tabular}
\end{flushleft}

\vspace{2em}
\noindent
All rights reserved. No part of this book may be reproduced, stored, or transmitted
in any form or by any means without the written permission of the author,
whether electronic, mechanical, by photocopying, recording, or otherwise.
\setlength{\parindent}{0pt}

\chapter*{Preface}
\addcontentsline{toc}{chapter}{Preface}

To many people, mathematics first appears as something abstract:
a world of formulas, symbols, and rules
whose connection to reality is not always immediately visible.
Yet this is precisely where its special strength lies.
Mathematics is not merely a technical tool,
but an inner framework
with which we order, understand, and, within limits, predict the world.

The natural sciences show again and again
that their laws are not arbitrary.
They are based on mathematical relationships
that we do not freely invent,
but gradually uncover.
Whether it is the number~$\pi$,
which has appeared in geometry and physics since antiquity,
the complex numbers,
which became the language of quantum mechanics,
or probability theory,
which gives precision to our understanding of chance --
everywhere we encounter the same experience:
mathematics captures real structures.

This book aims to make this dual role visible:
mathematics as an \textbf{abstraction of thought}
and at the same time as the \textbf{structure of nature}.
It follows how mathematical ideas arise,
how they expand,
and why they are able to grasp reality so deeply.
The focus is not only on computational methods or formulas,
but on the guiding question
of whether mathematics is merely invented
or whether it reveals an order
that is embedded in the world itself.

If this book is meant to lead to one insight in the end,
it is this:
mathematics is not separate from reality.
It is one of the clearest forms
in which reality can be recognized as structure.

\vspace{2em}

\begin{flushright}
	\textit{Dipl.-Ing.\,(FH) Christian Weilharter}\\
	\vspace{0.5em}
	Traunstein, 2026
\end{flushright}
\newpage
\noindent
\bigskip
\noindent
\textbf{Reading Note}

This book follows a continuous thread.
Concepts are introduced,
developed through examples,
and condensed at the end of each chapter in terms of their significance.
The boxes in the text mark central points
such as definitions, historical context, common misconceptions, or key statements.
They are intended as orientation aids,
but they do not replace the main text,
where the actual argument is developed.

\medskip
\noindent
For better orientation, the box types are distinguished by color:
note boxes provide guidance,
history boxes place concepts and developments in their historical context,
math boxes summarize formal content,
and didactic boxes highlight key conceptual points.